% Same-family working fragment, 25 October 2026. Root authored.
% Independent review and integration pending; shared B17 local lemmas separate.
\subsubsection{\(D_{17}\) with an actual short rank-one factor}
\label{app:D17I18}
Keep the actual roots and signed fork coordinates of
Appendix~\ref{app:D17ordinary18}. The original vector is
\(z=E(p-x-ye)\), with \(\norm x^2+y^2<1\). Its actual isolated
rank-one opposite is
\[
 c_I=\sqrt h\,p-\sqrt{1+h}\,e,\quad0<h<1/2,\qquad
 s^2={\sqrt{1+h}-1\over\sqrt{1+h}+1}.
\]
We have \(y\ge s\), \(x_1\ge t\), where
\[
 t={\sqrt{1+14s^2+s^4}-4s\over1-s^2},\qquad t+2.5s>.923
\]
by Lemma~\ref{lem:I923-real-barrier}. The same \(y\) coordinate
is counted once; the ordinary debit \(y^2>.1\) is unavailable.

Deleting the head retains actual \(D_{16}\). Its complete extension
list is isolation or simple three at long leaf two. The local nulls,
points and native-gap arguments from the preceding subsection retain
old \(I\) and both forks. Thus Lemmas~\ref{lem:I-finite-head-flat-reference}
and~\ref{lem:I-head-flat-reference} apply on any legally funded
reference constant first five, with mean \(<.361\) and individual
\(P_5\) coefficient \(>.146\). Their head-plus-\(I\) gain is
\(>.07665\), supported in first five plus \(I\), without assuming
strong five or any \(B\)-terminal conclusion. Finite isolated heads
have \(h>1/4\), \(y^2>.0557\); pure owns five/six are impossible.

\paragraph{Fresh tails before strong five.}
The rank-valid \(M=11\) inventory covers five through twelve, with
only nonadjacent pure far helpers weak. There is a strong index in
each of \(9/10,10/11,11/12\). The geometric \(D_4/D_3\) lists
from Appendix~\ref{app:D17ordinary18} precede its ordinary energy
debit. Their only extra weak rows are pure own-five/fork-three at
thirteen and pure own-five or finite-own-at-least-seven/end-three
at fourteen. Both forks and every foreign contact were retained;
dotted end-three owns were directly strong.

The actual left-\(H_4\) finite points are flat through \(i-3\),
have norm \(<i\) and demands \(>3.48,3.61\) at \(i=13,14\).
If five is strong, their preceding one-square totals before an
ordinary debit are \(>.9757\); adding \(>.07665\) excludes them.
If five is weak, the actual pressures \(x_1>.5965,S_6>2.01815\)
give the stronger bare totals \(>1.039,1.056\) on the positive
\(P_1,P_6,P_K\) references. No head gain is needed in that branch.

For every remaining \(.9<g<1\) end-three own, average the two
actual \(A_2\) cores, retaining their full cross \(8g^2/9\).
The demand is \(>3.388\), norm \(<14\), flat through twelve.
In the strong-five branch pay \(.16/.15\) at actual strong ten/eleven,
below the individual terminal prices. With
\(q_{10}>.823735,q_{11}>.8472225\), the bare totals are
\[
 .823735+.32(3.388-2.732)-.0256(4)=.931255,\quad
 .8472225+.30(3.388-2.886)-.0225(3)=.9303225.
\]
Adding the local head gain closes both. In the weak-five branch the
reference norms are \(>.8870,.9104\); the capped initial-\(I\)
graft, Lemma~\ref{lem:I-capped-initial-graft}, adds \(>.015\).
Its support first six plus \(I\) is disjoint from the tail residual.
Its \(P_6\) price is available, and the totals exceed
\(1.00952,1.0085\). This removes the tails before any later floor.
The pure-helper two-cap calculation extends nonadjacency through
fourteen: the actual inverse is \(2\min(i,j)\) and both retained
fork blocks remain positive. Its unchanged loss is
\(<7/16+1/(4\sqrt{10})<.61\). Thus thirteen or fourteen is strong.

\paragraph{Retained weak-five and weak-fourteen feet.}
Suppose five weak. For a justified \(x_1>a_0\), interpolate through
one, six and actual strong knots up to \(K=13,14\), omitting weak
ones. Put \(b_1=(f_6-a_0)/5\), \(C(a_0)=a_0^2+5b_1^2\).
The increment sums \(>.19304,.21155\), minus at most three/four
losses each \(<1/3200\), give
\(q_{13}>C+.1921025,q_{14}>C+.2103\). The needed capacities are
\[
\begin{array}{c|c|c|c|c}
a_0&C>&q_{13}>&P_1\text{ coefficient}>&P_6\text{ coefficient}>\\\hline
.5965&.7600299945&.9521324945&.312168&.08806\\
.643&.7916565045&.9837590045&.367968&.07876\\
.65&.7968668845&.9889693845&.376368&.07736
\end{array}
\]
Terminal thirteen/fourteen coefficients are \(>.14144,.1361\).
We never assign a terminal coefficient to an internal knot.

If fourteen is weak, its pure far-helper time is
\(A_{14}^2=28g^2+2g-1/2\). The finite/dotted feet use actual long
roots thirteen/fourteen/fifteen and the initial local roots.
After projection on \(P_1,P_6,P_{13}\), their initial pieces
\(H_3,H_2,H_1\) have respective \(P_1,P_6\) coefficients
\[
 (\sqrt2/10,\sqrt2/15),\quad
 (\sqrt2/5,\sqrt2/20),\quad(1/\sqrt2,0).
\]
The residual norm bounds for own five, seven/eight, finite at least
nine, and dotted are \(.824,1.105,1,1.28\). Only the last two
residuals are confined to fourteen/fifteen; initial crosses in the
other rows must be retained. The actual time/radical bounds are
\[
\begin{array}{c|c|c}
\text{own}&A_{14}>&\text{radical}<\\\hline
5&4.4095&.111\\7&4.9017&.249\\8&5.0242&.322
\end{array}
\]
For finite own at least nine the deficit increases with \(g\ge.939\).
At that endpoint \(A>5.1055\), the other time radical \(>5.116\)
and subtracted radical \(<.344\) give demand \(>4.7838\).
Dotted weakness implies \(g<10/9\): the squared ratio there is
\(83521/34020>(39/25)^2\), above the exact target, which is
\(<39/25\). The deficit decreases on this whole interval, since
\(A'<5.4\) and the radical derivative exceeds one. At \(10/9\),
\(A^2=5879/162>6.023^2\), radical \(<.436\), so demand \(>5.587\).
We use \(\sqrt{26}-.61<4.4891\) throughout.

Every finite-own weak-five foot gives \(x_1>.65\). Fourteen strong
then gives \(q_{14}(.65)>1\). Otherwise the weak-fourteen gains
are \(>.06160,.04163,.03203\) for owns five/seven/eight. For finite
at least nine the deficit is
\(>4.7838-.939(4.4891)-.460=.1085351\), and weight \(.09\)
gives gain \(>.011436318\), total \(>1.0004057\). The dotted
deficit \(>.1391\) gives \(>.01467\). Their prices fit the table.
Hence weak five is dotted, \(1<a<6/5\). Retain its full point
\[
 U_5=\sqrt2aP_5-a^{-1}n_5+\Omega_1/2,\quad
 T_5=\sqrt{10a^2+2a-1/2}-\sqrt{1-a^{-2}}.
\]
Its projection is \(mP_6+lP_1\), with
\(m=4\sqrt2a/5<1.35766,l=\sqrt2a/5+1/\sqrt2<1.04652\).
The residual has
\[
 Y_5^2=8a^2/5-2+a^{-2}<1,\quad Y_5\cdot P_1=Y_5\cdot P_6=0,
 \quad(Y_5)_2<.339412.
\]
On the \(.5965\) reference, its deficit \(>.1203\), at weight
\(.05\), gives gain \(>.00953\) and \(P_1/P_6\) debits
\(<.052326/.067883\).

A nonisolated head forces \(x_1>.643\). If fourteen strong,
\(q_{14}(.643)>1\). If weak, on \(q_{13}(.643)\) the retained
five foot gives \(>.00466\), while the fourteen finite-at-least-nine
and dotted gains are \(>.0126789498,.01557\), with zero cross.
For owns five/seven/eight the full cross losses \(.00144,.00108\)
leave gains \(>.06034,.04091,.03131\). Combined \(P_6\) price
is \(<.076369<.07876\), so all contradict one. A finite isolated
head instead has \(y^2>.0557\), already too much beside
\(q_{13}(.5965)\). Only isolated dotted heads remain under weak five.

On the \(.5965\) reference the full weak-fourteen accounting is
\[
\begin{array}{c|c|c|c}
\text{own}&\text{deficit lower}&\text{weight}&\text{net gain}>\\\hline
5&.387329&.09&.06160482\\7&.2960209&.09&.043253262\\
8&.24272051&.09&.0336591918\\
\ge9\text{ finite}&.1465351&.099&.0192\\
\text{dotted}&.1771&.09&.02151
\end{array}
\]
Here \(Y_5\cdot Y_{14}=2a/15,a/10,0\) for five, seven/eight,
and the last two rows. The first two totals already exceed one.
For own eight add an initial-\(I\) graft capped at \(.08\).
Its unclipped gain is \(>.015\) by weighted Cauchy. If clipped,
\(t>.5965+1/15\); since \(t(.11)<.65\), we have \(s<.11\) and
\[
 s^2+.16(t-.5965)-(5/6)(.08)^2
 >(s-.2)^2+.00690666>.015.
\]
Its residual is orthogonal to both feet because \(P_1,P_6\) are
retained. The full \(P_6\) price \(<.087582<.08806\) includes
the extra \(.013334\). Thus weak fourteen leaves only the last two
rows, each with gain \(>.0192\), \(P_1\) price \(<.070004\),
\(P_{13}\) price \(<.141422\), no \(P_6\) price and a residual
supported only on fourteen/fifteen.

\paragraph{Dotted heads on the retained weak-five reference.}
For a pure dotted head put \(u=g^{-2}\), \(h_0=u/[4(2-u)]\),
\(s_0=s(h_0)\). The actual native gap gives \(h\ge h_0\). The
joint point of Lemma~\ref{lem:I-pure-dotted-joint-point} supplies
\[
 ax_1+v_0y\ge F=1-d,\quad a^2=1-u/2,\quad
 v_0^2={u\over8-3u},\quad d^2=u/2-v_0^2.
\]
Besides its stated all-\(h\) domain \(u\ge.2\), it is valid for
\(.1\le u\le.2,h\le.2\). Indeed \(u(8-3u)\ge.77>2/3\)
and \(h/(1+h)\le1/6\), so
\(D_0=u/2-v_0\sqrt{h/(1+h)}>0\). The actual construction is
\[
 p+\rho c+\sigma c_I=D_0^{-1}(p-ae_1-v_0e),\quad
 \rho={a\over D_0\sqrt2g}>0,\quad
 \sigma={v_0\over D_0\sqrt{1+h}}>0.
\]
The cap product is zero, the \(c_I\) product has sign
\(v_0-\sqrt{h/(1+h)}\le0\), old \(n_1,e\) products are negative
and every unknown-wall product is negative. Both actual D forks
therefore remain satisfied, for every mutual head/\(I\) label.

Put \(X=(F-v_0s_0)/a\). Simplification gives
\[
 X=\sqrt{{1-\sqrt{3u/8}\over1+\sqrt{3u/8}}},\quad
 (1-X^2)(1+s_0^2)=4\sqrt3s_0X,\qquad X+2s_0>.91.
\]
For completeness, \(0<s_0<\sqrt5-2<.237\). At
\(X_0=.91-2s_0>0\), the left-minus-right expression is greater than
\(.1719-2.66467s_0+10.0283s_0^2+3.64s_0^3-4s_0^4\).
Writing \(w=s_0-.13\), this is exactly
\(.00182581+.092084w+11.0423w^2+1.56w^3-4w^4\), bounded below
by \(.00182581-.092084|w|+10.7719w^2>0\). The discriminant
is negative. The original expression decreases strictly with
positive \(X<1\), proving the full barrier without sampling.

Let \(W_0\) be the reference plus \(.05Y_5\) and any remaining
fourteen foot. Its first coordinate is still \(a_0=.5965\), and
the feet are orthogonal to \(P_1,P_6\). Set
\(\delta=\min((X-a_0)_+,.1)\), \(\kappa=6\delta/5\).
The joint paid-reference lemma uses nonnegative multipliers
\(\eta=\kappa/a\), \(\tau=s_0-\eta v_0\) and gives gain
\[
 s_0^2+2\kappa(X-a_0)-(5/6)\kappa^2.
\]
Here \(s_0/v_0>1/2\), \(a\ge1/\sqrt2\), \(\kappa\le.12\)
ensure \(\tau\ge0\). The \(P_6\) price \(\le.02\) makes the
full cost \(<.087883<.08806\); the other prices still fit.
If unclipped, the barrier and weighted Cauchy give gain
\(>(.91-.5965)^2/(4+5/6)>.02033425\). If clipped,
\(X>.6965>X(1/3)\) gives \(s_0<.111\); the gain exceeds
\((s_0-.24)^2+.00564>.022281\).
For \(u\le.1\), the actual head edge gives \(x_1>.7745\),
and an initial graft without an \(I\) charge gains \(>.03072\).
For \(.1\le u\le.2,h>.2\), \(I\) alone costs \(>.0455\).
Thus all pure dotted parameters are covered. The least totals are
\[
 .9703299945+.00953+.02033425>1.0001942,\quad
 .9521324945+.00953+.0192+.02033425>1.0011967.
\]

For a positive foreign contact put \(Q=\beta^2\ge1/4\),
\(v=(1-Q)/g^2\in(0,.75)\). Actual acuteness gives
\(h\ge Qv/[2(1-Q)-v]\ge v/(6-4v)\). The positive core
\(c+\beta e+((1-Q)/g)n_1\) clears the whole old \(I\), with
\[
 x\cdot Z\ge H(v)=\sqrt{1-v/2}-\sqrt{v(1-v)/2},\quad
 Z=(1-v/2)e_1+(v/2)e_2.
\]
Its projection is \((1-3v/5)P_1+(v/10)P_6\), residual norm
\(v^2/5\). The full retained initial cross is
\(W_0\cdot Z\le.5965-.14759v\): use
\(b_1<.284332,.05(Y_5)_2<.016971\). The remaining fourteen
foot has zero cross. Pay \(\nu=.18\) for \(v\le.2\),
\(.21\) for \(.2\le v\le.4\), and nothing thereafter. The gain is
\(2\nu[H-.5965+.14759v]-\nu^2v^2/5\). Adding \(I\) once gives
\[
\begin{array}{c|c|c|c}
v\text{ interval}&H>&I\text{ energy}>&\text{total gain}>\\\hline
(0,.2]&.6658&0&.0246888\\[.2,.25]&.6292&0&.02557275\\
[.25,.3]&.5979&.0121&.02739115\\[.3,.35]&.5709&.01515&.02191389\\
[.35,.375]&.5590&.0182&.0229054175\\
[.375,.4]&.5480&.0200&.021464225\\
[.4,.75)&\text{unused}&.0217&.0217
\end{array}
\]
Positive squares at right endpoints and monotonicity on \((0,.5)\)
prove the \(H\) bounds. The left endpoints in the lower bound for
\(h\) give the \(I\) bounds. The largest full \(P_1\) price is
\(<.052326+.070004+.1848=.30713<.312168\); the \(P_6\) price is
\(<.067883+.0084<.08806\). The terminal price is unchanged.
The uniform gain \(>.0214\) contradicts one in both fourteen
branches. Every weak-five head is excluded. Five is therefore strong.

The pure-helper coverage through fourteen is now available. If centre
fifteen were strong, its reference would have \(q_{15}>.9240775\),
first-five mean \(<.361\) and \(P_5\) coefficient \(>.146\).
The local head-plus-\(I\) gain would give \(>1.0007275\).
Thus centre fifteen is weak.

\paragraph{The complete centre-fifteen list.}
Deleting fifteen retains actual \(A_{14}+2A_1+I\). Its cap has form
\[
 c=Ap-g\Omega_{15}-\eta\Omega_j-\delta\Omega_{16}
                       -\epsilon\Omega_{17}-\beta e.
\]
There is at most one left contact. Put \(Q=\beta^2\); the exact data are
\[
 q_L=2\eta^2j(15-j)/15,\quad
 \Delta=1-q_L-\delta^2-\epsilon^2-Q>0,\quad
 \sigma=g+\eta j/15+(\delta+\epsilon)/2,\quad
 C_{15}=\sqrt{30}\sigma/\sqrt\Delta.
\]
The sufficient strong ratio is \(<T(11/2)<1.554\).
With no left contact, finite classification permits isolation, one
arbitrary finite fork label, or pairs \((3,3),(3,4),(3,5)\).
Without forks the left possibilities are endpoint three/four at
one/fourteen or next-to-end three at two/thirteen. With forks only
left endpoint-three with one fork-three/four or both forks-three
is finite. Two left contacts make a cycle; other branches or internal
multiple bonds are absent from the full finite classification at this rank.
Positive \(I\) contacts in a single-fork row are precisely
\((m,I)=(3,3),(3,4),(3,5),(4,3),(5,3)\); with two forks only
\((3,3,I3)\) is possible. Left endpoint-three may carry \(I3/I4\),
and endpoint-three/fork-three may carry \(I3\); other left bridges
are sealed. We retain the whole foreign alphabet.

For isolation, own zero has negative time; three/four/six/parallel
have the actual star/path/null principals. On \(0<u<1\), the cleared
core has pressure \(\sqrt{30-u}-\sqrt{u(1-u)}\), loss \(<.6\).
At \(u=1\) it is a forbidden feasible future null. On \(u>1\),
the actual finite projection gives \(\sqrt{29}-\sqrt{u-1}\).
Foreign cost forces \(u<75/64\); pure native at least seven has
\(u<100/81\), with losses \(<.515,.6\). Pure own five uses the
actual left-\(H_4\) projection, whose demand is
\(>5.3/1.04>\sqrt{30}-.61\). Every isolated centre is strong.

For left-only own zero, far-three needs \(1/2<Q<8/15\), containing
no native foreign cost. Far-next-three and far-four have zero time
at \(Q=0\); their admissible nonzero costs are too large. The near
three/next-three/four rows have actual affine \(D_5,D_6,B_4\)
principals. Positive owns three/four/six and parallel have the
actual cycle/star/path/null dispositions. Every other near/next/four
ratio is strong. A positive-foreign far-three has ratio
\(>(.809+1/30)/\sqrt{17/60}>1.57\).
For left plus fork-three, \(\Delta\le17/60\): own zero/three has
a cycle/triangle; larger owns are strong. With fork-four or both
forks-three, \(\Delta\le1/30\) and strength holds even at own zero.
Only the pure far helper can remain among left rows.

For own zero with one fork, fork labels three/four/six have the
actual affine \(E_8,B_4\) and \(6,3\) principals. For finite
fork cosine \(b\ge\cos(\pi/7)>.9\), whiteness forces \(Q=0\).
Set \(u=b^{-2}\in(1,100/81)\). The actual positive future core
\(c+b^{-1}n_f+2(b^{-1}-b)n_{15}\), divided by \(b/\sqrt2\),
has signed displacement \((1^{14},3-2u,u-1,\pm(1-u))\), with
\[
 G=\sqrt{17-2u}-\sqrt{-8+14u-6u^2},\quad R=25-16u+6u^2.
\]
Both fork signs are retained. Complete interval bounds are
\[
\begin{array}{c|c|c}
u&G>&R\le\\\hline
[1,1.1]&3.472&15\\[1.1,1.15]&3.4277&14.66\\
[1.15,1.2]&3.41265&14.535\\[1.2,100/81]&3.4119&14.44
\end{array}
\]
Use the maximum second radical \(1/\sqrt6\), positive squares and
decreasing \(R\). Pay \(.14/.13\) at actual terminal thirteen/
fourteen, with residual norm \(\le R-K\). Adding \(>.07665\)
makes all eight totals exceed \(1.00135\).

For fork five put \(b=\phi/2\), \(\alpha=b/\sqrt2\),
\(\delta_5=1-3b^2/2\). The actual finite \(H_4\) consists of cap,
contacted fork, centre and other fork. Its positive correction is
\(b(3n_f/2+n_{15}+n_{\rm other}/2)\). The corrected vector is
spacelike; the point is the normalized positive finite projection
of \(p\). Its pure displacement is \(\alpha P_{14}\).
For \(Q=0\), demand \(>3.49\), norm14, and the local head gain
close both references. The genuinely permitted \(Q=1/4\) retains
\(I\) in this same point:
\[
 x\cdot P_{14}+dy\ge G>3.61,\qquad
 d=1/(2\alpha)<.875,\quad d^2<.765,
 \quad G={\sqrt{5+6\phi}\over1+\sqrt{(13-8\phi)/(18-2\phi)}}.
\]
At \(K=13,14\), weights \(.14,.13\) give bare mixed gains
\(>.0874216,.0649831\). Every dotted/foreign head has pure gain
\(>.04\), and a nonisolated head has \(>.0994\), already enough.
A pure finite head has pure gain \(>.0186\) and \(s>.236\).
Append only the shifted \(I\) gain \((s-\mu d)^2\), exceeding
\(.1135^2,.12225^2\). This replaces the existing coefficient
\(\mu d\) by \(s\), using a nonnegative additional multiplier;
it does not add a second \(y^2\). The least totals exceed
\(1.00754385,1.0056781625\). Head and pure tail residuals are
orthogonal, and the displayed mixed norm includes the \(I\) coordinate.

For positive own with one fork, own-three/fork-three is an affine
triangle. Own-three/fork-at-least-six, or fork four/five with foreign
contact, is directly strong. With fork four/five and \(Q=0\), the
same cap gives the actual opposite of the contacted fork. The positive
finite \(D_{17}\) replacement is future feasible against every
original wall; the omitted fork is strict. Its actual cosh is
\(15/2+(17/2)b>13\). Thus the signed direction
\(P_{16}\pm e_{17}\), of norm17, has demand
\(>\sqrt{102/7}>3.817\). Either funded reference and the local
head gain close it. Own four/six/parallel have the actual star/path/
null principals. For larger owns, every fork label at least four is
directly strong.

Only fork-three remains. Own five with \(Q\ge1/2\) is strong.
At \(Q=0,1/4\), use the actual left-\(H_4\) point, or respectively
the cleared future core
\[
 c+g(3n_{15}/2+n_{14}+n_{13}/2)+e/2,\qquad g=\phi/2.
\]
It gives a pure displacement flat through twelve, norm \(<15\),
demand \(>3.52\). Weights \(.15/.14\) at actual strong eleven/
twelve, plus \(>.07665\), give totals \(>1.0240725,1.0229116\).
For every larger finite/dotted own, foreign contact is directly strong,
as is \(g\ge6/5\). On the full remaining \(Q=0,.9<g<6/5\)
interval use the actual core
\[
 \xi=c+g^{-1}n_{15}+(2g)^{-1}n_{14},\quad
 k=1-3/(4g^2)>0,\quad \alpha=\sqrt2(g+1/4).
\]
Its cap and fourteen products vanish; its fifteen and both fork
products are negative. Its square is \(-k\), time positive, and
every unknown-wall product is nonpositive by positive actual-wall
coefficients. With \(A^2=15\alpha^2-3/4\),
\(A/\alpha>3.836\), \(k/\alpha^2\le575/5046<.34^2\), the demand
is \(G>3.496>3.49\), full norm \(<15\), flat through thirteen.
The quotient derivative reduces to \(2g^3-3g-3/8<0\) throughout.
At thirteen pay \(.14\), residual norm \(<2\). At fourteen the
fourteenth coordinate is \(1-a\),
\(a=[4g(g+1/4)]^{-1}<1/4\). Its projection is
\((1-a/14)P_{14}\), reference pairing \(\le f_{14}\), residual
norm \(<1+2a<1.5\). Weight \(.13\) is affordable at that actual
terminal; the least total is \(>1.0051616\). No internal terminal
coefficient is assumed.

For two forks the finite pairs are \((3,3),(3,4),(3,5)\).
Own three has a triangle; owns at least four are directly strong.
Own-zero pair-three-three has the actual affine square, including its
foreign case. The other pairs are foreign-sealed. Pair-three-four
gives an actual reciprocal finite \(B_{17}\) endpoint, with
cosh \(15/2+(17/4)\sqrt2>13\), hence the same signed demand above.
Pair-three-five uses the actual finite \(H_3\) correction
\(n_{f_1}/2+(\phi/2)n_{f_2}\). Its pure displacement is
\(\alpha P_{15}\), \(\alpha=(1+\phi)/(2\sqrt2)\), and the
positive finite residual satisfies \(q=\delta_H/\alpha^2=2\phi^{-6}<1/8\).
Its pressure
\(\sqrt{15-q}/[1+\sqrt{q/15}]>3.53\), norm15, closes either
funded reference with the local head gain. This exhausts the list.

\paragraph{The final pure helper and the actual-wall interface.}
The weak centre must therefore be the pure far-three helper at fifteen,
\(A^2=30g^2+2g-1/2\). Nonadjacency with fourteen follows from the
same actual prefix Schur calculation, retaining both positive fork
blocks. Hence fourteen is strong, \(q_{14}>.90715\). Nonisolated
and finite isolated heads already contradict this reference by their
local gains. For dotted heads use
Lemma~\ref{lem:I-helper15-head-coupling}, checking its local geometric
interface as follows.

Its finite/dotted helper vectors are
\(g\sqrt2P_{15}-gn_{15}+H_k\) and
\(g\sqrt2P_{15}-g^{-1}n_{15}+H_1\). They use only actual long
\(n_{15},n_1,n_2\); \(P_{16}\) in their coordinates is a spatial
sum, not a nonexistent long D dual. Both actual D forks have
nonpositive product with every positive correction. Their Gram data,
fresh times, initial crosses and individual \(P_5/P_{14}\) prices
are literally those of the local kernel, without requiring a B short
normal or a sign for \(x_{17}\).

For pure dotted head \(u\le.2\), the kernel bypasses the tail payment
and gives gain \(>.11088\). For \(.2\le u<1\), its net tail gain
is \(>.0421\), including the whole initial residual, and its
head-plus-\(I\) gain is \(>.0515\). Every positive foreign head
contact instead gives \(>.0521\), with the adverse initial cross
and the one \(I\) cost included. The least original total is
\[
                   .90715+.0421+.0515=1.00075>1.
\]
The \(P_5\) debits \(<.09873,.09273\) fit \(>.146\), and
\(P_{14}\) debit \(<.1273\) fits \(>.1361\). Every actual centre,
head, finite foreign label and own range is covered on the same
original vector. This excludes \(D_{17}+I\) in dimension eighteen,
without arithmeticity or an ordinary rank-one debit.
